% sections/comparison.tex: Section 7: the estimates against prior work. Table from make_comparison_table.R.
\section{Comparison with Prior Estimates}\label{sec:comparison}

Table \ref{tab:comparison} in Appendix \ref{appendix:figures} places the estimates from this paper alongside prior discount-rate estimates. The rows are not all the same object, and the table says which each is: an exponential rate over a stated horizon, a present-payment premium, a bound, or a laboratory elicitation whose horizon is the experiment's.

The section makes three comparisons. First, the paper's estimate of the rate over the contract horizon, $0.11$ with a set from $0.09$ to $0.13$ from the within-borrower design, lies inside the set the rescheduling experiment accepts for distressed borrowers over four years, zero to $1.2$ at the five percent level; the experiment fixes the sign of the far-payment weight, and the bureau's contract fixes its level, because the auto loan's horizon is where the rate is measured most precisely (Section \ref{sec:mapping}). Second, the laboratory literature's median of 22 percent per year \citep{cohen_measuring_2020} and the field elicitations of 20 to 28 percent are elicited over horizons of weeks to a year. Over those horizons this paper's object is the present-payment premium, under which a dollar due next month weighs nine tenths of a dollar due now for a household at the default margin and a dollar due at the end of a five-year contract about a third, the constant-income level of the premium being the fifth that Indarte's pair measures; the step at the first payment is larger than the elicitations report, while the exponential rate beyond it is about half of them. The two literatures measure different segments of the same schedule, and the schedule has a large step at the first payment and a low rate thereafter. Third, against the asset-market calibrations, the paper's $0.11$ is more than five times the two percent social rate and near the zero-beta rate of the marginal investor; the present-payment premium, not the rate, is what most distinguishes borrowing households from the households whose preferences asset prices reveal.
